\documentclass{article}

\usepackage{amsmath,amssymb}
\usepackage{xurl}
\usepackage[colorlinks=true,linkcolor=blue,citecolor=blue,urlcolor=blue]{hyperref}
\setlength{\emergencystretch}{2em}

% Shared commands for notes in docs/paper-gaps/.
% Notation follows the LDT paper (references/ldt-paper/).

\newcommand{\Lean}{\textsc{Lean}}
\newcommand{\leanid}[1]{\textsf{#1}}
\newcommand{\ghissue}[1]{\href{https://github.com/LionSR/MIPStarRE/issues/#1}{GitHub issue \##1}}
\newcommand{\ghpr}[1]{\href{https://github.com/LionSR/MIPStarRE/pull/#1}{GitHub PR \##1}}

% --- Paper-compatible notation ---
\newcommand{\F}{\mathbb{F}}
\newcommand{\N}{\mathbb{N}}
\newcommand{\C}{\mathbb{C}}
\newcommand{\tr}{\operatorname{tr}}
\newcommand{\ot}{\otimes}
\newcommand{\E}{\mathop{\bf E\/}}
\newcommand{\bx}{\mathbf{x}}
\newcommand{\by}{\mathbf{y}}
\newcommand{\bu}{\mathbf{u}}
\newcommand{\bv}{\mathbf{v}}

% Approximation relation (paper uses \approx_\eps and \simeq_\zeta)
\newcommand{\approxeps}{\approx_{\varepsilon}}
\newcommand{\simgeq}{\simeq_{\zeta}}

% Polynomial spaces (paper uses \polyfunc, \polysub, \polymeas)
\newcommand{\polyfunc}[3]{\operatorname{Poly}(#1,#2,#3)}
\newcommand{\polysub}[3]{\operatorname{PolySub}(#1,#2,#3)}
\newcommand{\polymeas}[3]{\operatorname{PolyMeas}(#1,#2,#3)}


\title{The Constant in the Formalized Orthonormalization Theorem}
\author{}
\date{2026-05-11}

\begin{document}
\maketitle

\section{At a glance}

\begin{itemize}
  \item \textbf{Difficulty.} Medium.  The discrepancy is a scalar loss, but
  recovering the paper constant requires aligning the local repair lemma with
  the bipartite self-consistency hypothesis used in the source proof.
  \item \textbf{Estimated weight.} One reformulated repair theorem, with a
  moderate amount of project-local scalar and transport work.
  \item \textbf{Mathlib/project split.} Mostly project-local; Mathlib supplies
  the elementary real-power inequalities and matrix order infrastructure.
  \item \textbf{Key inputs.} The Section~5 completion step, the conversion from
  consistency to almost-projectivity, and the local \(Q/X/\widehat X/P\) repair
  theorem.
\end{itemize}

\section{Key theorem forms}

Throughout, \(\ket{\psi}\) denotes a normalized bipartite state on
\(\mathcal H_A\otimes\mathcal H_A\), \(A=\{A_a\}_{a\in\mathcal A}\)
is a submeasurement on \(\mathcal H_A\) with strong self-consistency error
\(\zeta\ge 0\), and \(P=\{P_a\}_{a\in\mathcal A}\) is a projective
submeasurement on \(\mathcal H_A\).  The relation
\(\{X_a\}\approx_\delta\{Y_a\}\), in the tensor-placement used here, abbreviates
the state-dependent squared-distance estimate
\[
  \sum_a
  \bra{\psi}
    (X_a\otimes I-Y_a\otimes I)^\dagger
    (X_a\otimes I-Y_a\otimes I)
  \ket{\psi}
  \le \delta.
\]

The paper states the orthonormalization theorem with the estimate
\[
  A_a \otimes I \approx_{100\zeta^{1/4}} P_a \otimes I.
\]
The present formal theorem proves the same conclusion with the estimate
\[
  A_a \otimes I \approx_{120\zeta^{1/4}} P_a \otimes I.
\]
The input-driven theorem, in which the stronger repair input is supplied
separately, still has the symbolic \(100\zeta^{1/4}\) envelope.%
\footnote{The source theorem is
\path{references/ldt-paper/orthonormalization.tex:67-75}
(\path{thm:orthonormalization}).  The relevant declarations are
\path{MIPStarRE.LDT.MakingMeasurementsProjective.orthonormalization},
\path{MIPStarRE.LDT.MakingMeasurementsProjective.orthonormalization_ofInput},
and \path{MIPStarRE.LDT.MakingMeasurementsProjective.leftLiftedProjectivizationRepairProducer}.
This note is attached to \ghissue{1032} and \ghpr{1480}.}

\section{The source assertion}

In the source proof of the orthonormalization lemma, the submeasurement \(A\)
is first completed to a measurement \(\widehat A\) on
\(\mathcal A\cup\{\bot\}\), by setting \(\widehat A_a=A_a\) for
\(a\in\mathcal A\) and
\(\widehat A_\bot=I-\sum_{a\in\mathcal A} A_a\).  The completed measurement
has bipartite strong self-consistency error \(2\zeta\).  The paper then applies
the main orthonormalization lemma directly at this error level.  The relevant
source passages are the main lemma statement at
\path{references/ldt-paper/orthonormalization.tex:282-293}
(\path{lem:orthonormalization-main-lemma}) and the completion-and-apply
argument at \path{references/ldt-paper/orthonormalization.tex:304-369}.
This gives
\[
  \widehat A_a \otimes I
    \approx_{84(2\zeta)^{1/4}}
  \widehat P_a \otimes I.
\]
The displayed estimate occurs at line \(358\), and the closing scalar
calculation \(84\cdot 2^{1/4}\simeq 99.89 \le 100\) occurs at line \(368\).
Thus the paper obtains the advertised \(100\zeta^{1/4}\) bound after
discarding the completion outcome.

\section{Comparison with the formalization}

The formalized unconditional theorem follows a slightly different route.  After
completion it converts the \(2\zeta\) self-consistency estimate into a
source-almost-projective estimate before invoking the locality-preserving
\(Q/X/\widehat X/P\) repair theorem.  This route is recorded by the named
Lean error function
\(\mathsf{orthonormalizationCompletionRouteError}(\zeta)=120\zeta^{1/4}\).
The conversion doubles the scalar:
\[
  \operatorname{consistencyToAlmostProjectiveError}(2\zeta)=4\zeta.
\]
The local repair theorem then gives
\[
  84(4\zeta)^{1/4}
    = 84\cdot 4^{1/4}\zeta^{1/4}
    \le 120\zeta^{1/4}.
\]
Thus the formalized theorem has the same mathematical shape as the paper
theorem, but its constant is weaker by the extra factor \(2^{1/4}\) introduced
by this conversion step.

\section{Conclusion}

The present statement is a documented weakening of the paper constant, not a
new hypothesis.  To recover the \(100\zeta^{1/4}\) estimate, the local repair
theorem should be reformulated so that it consumes the completed measurement's
bipartite strong self-consistency estimate directly, avoiding the intermediate
conversion to source almost-projectivity.

\bibliographystyle{alpha}
\bibliography{references}

\end{document}
